\documentclass[oneside]{ctexbook}
\usepackage{ctex}
\usepackage[utf8]{inputenc}
\usepackage{amsmath}
%\usepackage{amsthm}
\usepackage{ntheorem}
\usepackage{booktabs}
\usepackage{caption}
\usepackage{listings}
\usepackage[dvipsnames]{xcolor}
\usepackage{xeCJK}
\usepackage{bm}
\usepackage{fancyhdr}
\usepackage{graphicx}
\usepackage{amssymb}
\usepackage{mathrsfs}
\usepackage{titlesec, blindtext, color}
\usepackage{arydshln}
\usepackage{hyperref} 
\usepackage[OT1]{fontenc}
\usepackage{geometry}
\usepackage{comment}
\usepackage{extarrows}
\usepackage{inconsolata}
\usepackage{yhmath}
\usepackage{enumitem}
\usepackage[titletoc]{appendix}
\usepackage{verbatim}

\geometry{a4paper,scale=0.8}

\hypersetup{hidelinks,
	colorlinks=true,
	allcolors=black,
	pdfstartview=Fit,
	breaklinks=true
}

\newcounter{sectionCnt} 


\ctexset{
	 chapter = {
		name = {第,章},
		format = \linespread{1.0}\zihao{1}\bfseries\sffamily\centering,
		nameformat = {\fontsize{30pt}{0pt}\centering\vspace{15pt}},
		titleformat = {\centering},
		number = \chinese{chapter},
		numberformat = \sffamily,
		aftername = ,
		beforeskip = {7pt},
		afterskip = {18pt},
		aftername = {\\},
		pagestyle = plain,
	},
	section={
		%format用于设置章节标题全局格式，作用域为标题和编号
		%字号为小三，字体为黑体，左对齐
		%+号表示在原有格式下附加格式命令
		format+ = \zihao{2} \heiti \raggedright\centering,
		%name用于设置章节编号前后的词语
		%前、后词语用英文状态下,分开
		%如果没有前或后词语可以不填
		name = {,.\,},
		%number用于设置章节编号数字输出格式
		%\chinese输出section编号为中文，\arabic输出为阿拉伯数字
		%number = \Roman{sectionCnt},
		number = \chinese{sectionCnt},
		%beforeskip用于设置章节标题前的垂直间距
		%ex为当前字号下字母x的高度
		%基础高度为1.0ex，可以伸展到1.2ex，也可以收缩到0.8ex
		beforeskip = 1.0ex plus 0.2ex minus .2ex,
		%afterskip用于设置章节标题后的垂直间距
		afterskip = 1.0ex plus 0.2ex minus .2ex,
		%aftername用于控制编号和标题之间的格式
		%\hspace用于增加水平间距
		aftername = \hspace{0pt}
	},
	subsection={
		format+ = \zihao{-2} \kaishu  \raggedright,
		%仅输出subsection编号且为中文
		number = \roman{subsection},
		name = {,.},
		beforeskip = 1.0ex plus 0.2ex minus .2ex,
		afterskip = 1.0ex plus 0.2ex minus .2ex,
		aftername = \hspace{0pt}
	},
	subsubsection={
		%设置对齐方式为居中对齐
		format+ = \zihao{-3} \fangsong,
		%仅输出subsubsection编号，格式为阿拉伯数字，打字机字体
		number = \ttfamily\arabic{subsubsection},
		name = {,.},
		%beforeskip = 1.0ex plus 0.2ex minus .2ex,
		afterskip = 1.0ex plus 0.2ex minus .2ex,
		aftername = \hspace{0pt}
	}
}

\newenvironment{Appendices}[1]
{

	\ctexset{
		chapter = {
			name = {#1},
			format = \linespread{1.0}\zihao{1}\bfseries\sffamily\centering,
			nameformat = {\fontsize{30pt}{0pt}\centering\vspace{15pt}},
			titleformat = {\centering},
			%number = \chinese{chapter},
			number = ,
			numberformat = \sffamily,
			aftername = ,
			beforeskip = {7pt},
			afterskip = {18pt},
			aftername = {\\},
			pagestyle = plain,
		},
		section={
			%format用于设置章节标题全局格式，作用域为标题和编号
			%字号为小三，字体为黑体，左对齐
			%+号表示在原有格式下附加格式命令
			format+ = \zihao{-1} \heiti \raggedright\centering,
			%name用于设置章节编号前后的词语
			%前、后词语用英文状态下,分开
			%如果没有前或后词语可以不填
			name = {,.\,},
			%number用于设置章节编号数字输出格式
			%\chinese输出section编号为中文，\arabic输出为阿拉伯数字
			number = \Roman{section},
			%beforeskip用于设置章节标题前的垂直间距
			%ex为当前字号下字母x的高度
			%基础高度为1.0ex，可以伸展到1.2ex，也可以收缩到0.8ex
			beforeskip = 1.0ex plus 0.2ex minus .2ex,
			%afterskip用于设置章节标题后的垂直间距
			afterskip = 1.0ex plus 0.2ex minus .2ex,
			%aftername用于控制编号和标题之间的格式
			%\hspace用于增加水平间距
			aftername = \hspace{0pt}
		},
		subsection={
			format+ = \zihao{-2} \kaishu  \raggedright,
			%仅输出subsection编号且为中文
			number = \roman{subsection},
			name = {,.},
			beforeskip = 1.0ex plus 0.2ex minus .2ex,
			afterskip = 1.0ex plus 0.2ex minus .2ex,
			aftername = \hspace{0pt}
		},
		subsubsection={
			%设置对齐方式为居中对齐
			format+ = \zihao{-3} \fangsong,
			%仅输出subsubsection编号，格式为阿拉伯数字，打字机字体
			number = \ttfamily\arabic{subsubsection},
			name = {,.},
			%beforeskip = 1.0ex plus 0.2ex minus .2ex,
			afterskip = 1.0ex plus 0.2ex minus .2ex,
			aftername = \hspace{0pt}
		}
	}
}

\fancyhf{}
\fancyhead[L]{
	\begin{minipage}[c]{0.06\textwidth}
		\hyperlink{Index}{\includegraphics[height=7.5mm]{1.jpg}}
	\end{minipage}
	\begin{minipage}[c]{0.4\textwidth}
	线代冲刺笔记
	\end{minipage}}
\fancyhead[R]{
    \begin{minipage}[r]{0.1\textwidth}
	\href{https://qifengggg.github.io/}{奇峰}
    \end{minipage}
}
\cfoot{\hyperlink{Index}{\thepage}}

\pagestyle{fancy}

\fancypagestyle{plain}{
  \fancyhead{}
  \renewcommand{\headrulewidth}{0pt}
}
%\renewcommand{\qedsymbol}{$\blacksquare$}

\renewcommand{\labelenumi}{\roman{enumi}.}
\renewcommand{\labelitemii}{$ \circ $ }
\newcommand{\dis}{\displaystyle}
\newcommand{\Attention}[1]{\textcolor{red}{\bf{#1}}}
\newcommand{\red}{\textcolor{red}}
\newcommand{\sssubsection}[1]{\noindent\textbf{#1}}
\newcommand{\nextline}{\vspace{5pt}\newline}
\newcommand{\dsqrt}{\dis\sqrt}
\newcommand{\pprime}{{\prime\prime}}
\newcommand{\whichis}{\mathop=\limits^\triangle}
\newcommand{\R}{\mathbb R}
\newcommand{\qline}{\underline{~~~~~~~~~~}}

\newcommand{\Section}{\stepcounter{sectionCnt}\section}

\title{\includegraphics[scale=0.6]{1.jpg}\\ \textsf{线代冲刺笔记}}
\author{\href{https://qifengggg.github.io/}{奇峰}}
\date{之前}

\begin{document}

\theoremseparator{}
\newtheorem{def1}{定义}[section]
\newtheorem{theo1}{定理}[section]
\newtheorem{func1}{方法}[section]
\newtheorem{infer1}{推论}[section]
\newenvironment{proof}{\begin{itemize}\item \textbf{证明}
	
	}{\end{itemize}}
%\newenvironment{<环境名称>}[<参数个数>][<首参数默认值>]{<环境前定义>}
%							{<环境后定义>}


\newenvironment{Def}[1][\quad{}]{\begin{def1}\textbf{#1}}{\end{def1}}
\newenvironment{Theo}[1][\quad{}]{\begin{theo1}\textbf{#1}}{\end{theo1}}
\newenvironment{Func}[1][\quad{}]{\begin{func1}\textbf{#1}}{\end{func1}}
\newenvironment{Infer}[1][\quad{}]{\begin{infer1}\textbf{#1}}{\end{infer1}}
\newenvironment{Field}[1][\quad{}]{\noindent\newline\textbf{#1}}{}

\setlength{\parskip}{3pt}
\setenumerate[1]{topsep=2pt}
\setitemize[1]{topsep=2pt}

\frontmatter
\maketitle
\hypertarget{Index}{}
\tableofcontents

\mainmatter

\chapter{二重积分}

\section{交换积分次序}

$ \begin{cases}
    \textrm{直角 - 穿针法}\\
    \textrm{极坐标}\begin{cases}
        xy\textrm{坐标系 - 作圆弧}\\ 
        \theta r\textrm{坐标系 - 穿针法(当成直角坐标)}
    \end{cases}
\end{cases} $ 

二重变限积分求极限时，
$ \begin{cases}
    \textrm{洛必达结合交换积分次序}\\\textrm{二重积分中值定理(非零因子)}
\end{cases} $ 

\sssubsection{例题T6.2.2}

\begin{equation*}
    \begin{aligned}
        {\displaystyle\lim_{t\rightarrow 0^+}}\dfrac{1}{t^3}\int_0^t\mathrm{d}x
        \int_{x^2}^{t^2}\arctan(\cos(2x+3\dsqrt{y}))\mathrm{d}y &= 
        {\displaystyle\lim_{t\rightarrow 0^+}}\dfrac{1}{t^3}\iint\limits_{D}
        \arctan(\cos(2x+3\dsqrt{y}))\mathrm{d}x\mathrm{d}y
        \\&= {\displaystyle\lim_{t\rightarrow 0^+}}\dfrac{1}{t^3}
        \arctan(\cos(2\xi+3\dsqrt{\eta}))\cdot S_D \,(\xi,\eta \in D)
        \\&= {\displaystyle\lim_{t\rightarrow 0^+}}\dfrac{1}{t^3}
        \arctan(\cos(2\xi+3\dsqrt{\eta}))\cdot\dfrac{2}{3}t^3
        \\&= \dfrac{\pi}{4}\cdot\dfrac{2}{3} = \dfrac{\pi}{6}.
    \end{aligned}
\end{equation*}

\sssubsection{例题T6.3.2}

\begin{equation*}
    \begin{aligned}
        \int_0^\frac{\pi}{2}\dfrac{1}{\sqrt{x}}\mathrm{d}x\int_{\sqrt{x}}^{\sqrt{\frac{\pi}{2}}}
        \dfrac{1}{1+(\tan y^2)^{\sqrt 2}}\mathrm{d}y 
        &= \int_0^{\sqrt{\frac{\pi}{2}}}\dfrac{1}{1+(\tan y^2)^{\sqrt 2}}
        \mathrm{d}y\int_0^{y^2}\dfrac{1}{\sqrt x}\mathrm{d}x
        \\&= \int_0^{\sqrt{\frac{\pi}{2}}}\dfrac{2y}{1+(\tan y^2)^{\sqrt 2}}\mathrm{d}y
        \\&= \int_0^{\frac{\pi}{2}}\dfrac{\mathrm{d}t}{1+(\tan t)^{\sqrt 2}}
        \\&= \int_0^{\frac{\pi}{2}}\dfrac{(\cos t)^{\sqrt 2}\mathrm{d}t}{(\sin t)^{\sqrt 2}+(\cos t)^{\sqrt 2}} 
        \\\textrm{区间再现}&=\dfrac{1}{2}\cdot\dfrac{\pi}{2} = \dfrac{\pi}{4}.
    \end{aligned}
\end{equation*}

\sssubsection{例题T6.3.4}

\begin{equation*}
    \begin{aligned}
        \int_0^1\dfrac{x^b-x^a}{\ln x}\mathrm{d}x &= 
        \int_0^1 \mathrm{d}x\int_a^b x^y\mathrm{d}y \\ 
        &= \int_a^b \mathrm{d}y \int_0^1 x^y\mathrm{d}x
        \\&= \int_a^b \dfrac{\mathrm{d}y}{1+y} = \ln\dfrac{1+b}{1+a}.
    \end{aligned}
\end{equation*}

\sssubsection{例题T6.3.5}

此处以 $ \theta - r $ 坐标系交换积分次序。

\begin{equation*}
    \begin{aligned}
        \int_0^{2\pi}(\theta^2-1)\mathrm{d}\theta\int_{\frac{\theta}{2}}^\pi e^{r^2}\mathrm{d}r &=
        \int_0^\pi e^{r^2}\mathrm{d}r\int_0^{2r}(\theta^2-1)\mathrm{d}\theta \\ 
        &= \dfrac{8}{3}\int_0^\pi r^3e^{r^2}\mathrm{d}r - 2\int_0^\pi re^{r^2}\mathrm{d}r
        \\ &= \dfrac{1}{3}(4\pi^2-7)e^{\pi^2}+\dfrac{7}{3}.
    \end{aligned}
\end{equation*}

\section{二重积分的计算}

考虑
$ \begin{cases}
    \textrm{直角坐标系}\\
    \textrm{极坐标变换}\\
    \textrm{奇偶性}\\
    \textrm{轮换对称性}\\
    \textrm{形心公式}\\
    \textrm{平移变换}\\
    \textrm{雅可比行列式}
\end{cases} $ 

\sssubsection{被积函数为齐次函数 $ f(\dfrac{y}{x}) $ T6.6}

\begin{enumerate}
    \item[\textbf{例题}] 设有区域 $ D = \{(x,y)|x|+|y|\leq 1\}, $ 
    计算 $ \dis \iint\limits_D e^{\dfrac{|y|}{|x|+|y|}}\mathrm{d}x\mathrm{d}y. $ 
    \item[\textrm{\textbf{方法}}] 可以知道，
    \begin{equation*}
        \begin{aligned}
            \textrm{原积分} &= 4\iint\limits_{D_1}e^{\dfrac{y}{x+y}}\mathrm{d}x\mathrm{d}y \\ 
            &= 4\int_0^{\frac{\pi}{2}}\mathrm{d}\theta\int_0^{\frac{1}{\sin \theta + \cos \theta}}
            e^{\frac{\sin \theta}{\sin \theta + \cos \theta}}r\mathrm{d}r
            \\&= 2\int_0^{\frac{\pi}{2}} e^{\frac{\sin \theta}{\sin \theta + \cos \theta}} 
            \dfrac{1}{(\cos \theta + \sin \theta)^2}\mathrm{d}\theta
            \\&= 2\int_0^{\frac{\pi}{2}} e^{\frac{\sin \theta}{\sin \theta + \cos \theta}} 
            \mathrm{d}\dfrac{\sin \theta}{\sin \theta + \cos \theta}
            \\&= 2e^{\frac{\sin \theta}{\sin \theta + \cos \theta}}
            \Big|_0^\frac{\pi}{2} = 2(e-1).
        \end{aligned}
    \end{equation*}
\end{enumerate}

事实上，对于包含齐次函数 $ f(\dfrac{y}{x}) $ 的二重积分，都应该考虑利用极坐标变出 
$ \mathrm{d}\dfrac{A\sin\theta+B\cos\theta}{\sin\theta+\cos\theta}. $

对于 $ f(\dfrac{y^2}{x^2}),f(\dfrac{ax+by}{x+y}) $ 也是类似。

\sssubsection{极坐标压轴题 - T6.7}

\begin{enumerate}
    \item[\textbf{例题}] 设 $ f(x,y) $ 在区域 $ D = \{(x,y)|x^2+y^2\leq 1\} $ 
    上有一阶连续偏导数，在 $ D $ 的边界上 $ f(x,y) $ 恒为零，求
    $$
        {\displaystyle\lim_{\varepsilon\rightarrow 0^+}}\iint\limits_{\varepsilon^2\leq x^2+y^2\leq 1}
        \dfrac{xf'_x+yf'_y}{x^2+y^2}\mathrm{d}x\mathrm{d}y.
    $$
    \item[\textbf{方法}] 引入极坐标，则有
    \begin{equation*}
        \begin{aligned}
            &\dfrac{\partial f}{\partial \rho} = 
            \dfrac{\partial f}{\partial x}\dfrac{\partial x}{\partial \rho}
            + \dfrac{\partial f}{\partial y}\dfrac{\partial y}{\partial \rho}
            = f'_x\cos\theta + f'_y\sin\theta \\ 
            &\dfrac{\partial f}{\partial \theta} = 
            \dfrac{\partial f}{\partial x}\dfrac{\partial x}{\partial \theta}
            + \dfrac{\partial f}{\partial y}\dfrac{\partial y}{\partial \theta}
            = -f'_x(\rho \sin\theta)+f'_y(\rho \cos\theta)\\ 
            \Rightarrow & \rho \dfrac{\partial f}{\partial \rho} = 
            \rho\left[f'_x\cos\theta + f'_y\sin\theta\right]
            =xf'_x+yf'_y.
        \end{aligned}
    \end{equation*}
    那么，对于题设，有
    \begin{equation*}
        \begin{aligned}
            \textrm{原极限}&= {\displaystyle\lim_{\varepsilon\rightarrow 0^+}}
            \int_0^{2\pi}\mathrm{d}\theta\int_{\varepsilon}^1
            \dfrac{\partial f}{\partial \rho}\mathrm{d}\rho \\ 
            &= {\displaystyle\lim_{\varepsilon\rightarrow 0^+}}
            \int_0^{2\pi}f(\rho\cos\theta,\rho\sin\theta)\Big|_\varepsilon^1\mathrm{d}\theta\\
            &= -{\displaystyle\lim_{\varepsilon\rightarrow 0^+}}
            \int_0^{2\pi}f(\varepsilon\cos\theta,\varepsilon\sin\theta)\mathrm{d}\theta \\ 
            &= -\int_0^{2\pi}f(0,0)\mathrm{d}\theta = -2\pi f(0,0).
        \end{aligned}
    \end{equation*}
\end{enumerate}

\sssubsection{极坐标逆问题 - T6.8}

在极坐标不好做的情况下，将 $ (\rho,\theta) $ 转换回 $ (x,y) $ 即可。

\newpage

\sssubsection{奇偶性 - T6.9-10}

\begin{itemize}
    \item \textbf{奇偶性的推广}
    
    设 $ f(x,y) $ 在区域 $ D = D_1\bigcup D_2 $ 上连续，其中 $ D_1,D_2 $ 关于原点对称，则
    $$
        \iint\limits_Df(x,y)\mathrm{d}x\mathrm{d}y
        =\begin{cases}
            2\dis\iint\limits_{D_1}f(x,y)\mathrm{d}x\mathrm{d}y,\quad{} f(-x,-y) = f(x,y)\\
            0,\quad{} f(-x,-y) = -f(x,y)\\
        \end{cases}
    $$
\end{itemize}

\begin{enumerate}
    \item[\textbf{例题}]
    设区域 $ D $ 由曲线 $ y = x^3,y = x $ 围成。求
    $$
        \iint\limits_D\left[e^{x^2}+\sin(x+y)\right]\mathrm{d}x\mathrm{d}y.
    $$
    \item[\textbf{方法}] 
    由奇偶性，可以知道
    \begin{equation*}
        \begin{aligned}
            \textrm{原积分}&=2\int_0^1\mathrm{d}x\int_{x^3}^x e^{x^2}\mathrm{d}y\\ 
            &= 2\int_0^1e^{x^2}(x-x^3)\mathrm{d}x 
            \\&= \int_0^1(1-x^2)e^{x^2}\mathrm{d}x^2 = e-2.
        \end{aligned}
    \end{equation*}
\end{enumerate}

\begin{enumerate}
    \item[\textbf{例题}]
    对 $ D = \{(x,y)|x^2+y^2\leq 1, x+y\geq 0\}, $ 求
    $$
        \iint\limits_D\dfrac{1+2x^2+xy^2}{1+x^2+y^2}\mathrm{d}x\mathrm{d}y.
    $$
    \item[\textbf{方法}] 
    原积分等于 $ \dis \iint\limits_D\dfrac{1+2x^2}{1+x^2+y^2}\mathrm{d}x\mathrm{d}y+
    \iint\limits_D\dfrac{xy^2}{1+x^2+y^2}\mathrm{d}x\mathrm{d}y. $ 
    前一部分利用轮换对称性。对后一部分，以 $ y = x $ 将 $ D $ 
    划分为 $ D_1,D_2, $ 以 $ y = 0 $ 将 $ D_1 $ 划分为 $ D_{11},D_{12}. $ 
    那么，
    \begin{equation*}
        \begin{aligned}
            \textrm{后半部分积分} &= \iint\limits_{D_1}+\iint\limits_{D_2}
            \dfrac{xy^2}{1+x^2+y^2}\mathrm{d}x\mathrm{d}y\\ 
            &= 2\iint\limits_{D_12}\dfrac{xy^2}{1+x^2+y^2}\mathrm{d}x\mathrm{d}y\\ 
            &= 2\int_{0}^\frac{\pi}{4}\mathrm{d}\theta\int_0^1
            \dfrac{r^4\sin^2\theta\cos\theta}{1+r^2}\mathrm{d}r = \dfrac{1}{3\sqrt 2}
            \left(\dfrac{\pi}{4}-\dfrac{2}{3}\right).
        \end{aligned}
    \end{equation*}
\end{enumerate}

\sssubsection{轮换对称性的推广 - T6.11}

\begin{itemize}
    \item \textbf{轮换对称性的推广}
    
    $$
        \iint\limits_Df(x,y)\mathrm{d}\sigma = 
        \iint\limits_{D'}f(y,x)\mathrm{d}\sigma
    $$
    其中 $ D' = \{(y,x)|(x,y)\in D\}. $ 
\end{itemize}

\end{document}
